% !TeX spellcheck = en_GB

\documentclass[aspectratio=169,handout]{beamer}

\usepackage[utf8]{inputenc}
\usepackage[english]{babel}
%\usepackage{appendixnumberbeamer}

\usepackage{tikz}
\usetikzlibrary{shapes,arrows,positioning,calc}

\tikzstyle{startstop} = [ellipse, draw, fill=bg, text=structure, text width=3em, text badly centered, inner sep=0pt]
\tikzstyle{block} = [rectangle, draw, fill=bg, text=structure, text width=5em, text centered]
\tikzstyle{decision} = [diamond, aspect=3, draw, fill=bg, text=structure, text width=3em, text badly centered, inner sep=0pt]
\tikzstyle{input} = [trapezium, draw, fill=bg, text=structure, text width=5em, text centered, trapezium right angle=120]
\tikzstyle{output} = [trapezium, draw, fill=bg, text=structure, text width=5em, text centered, trapezium right angle=120]

\tikzstyle{line} = [draw, -latex']

%\usepackage{xcolor}
\usepackage{listings}

\lstdefinestyle{pseudo}{
	backgroundcolor=\color{bg},
	keywordstyle=\color{alert},
	numberstyle=\tiny\color{fg!50},
	numbers=left,
	%commentstyle=\color{mGreen},
	%stringstyle=\color{mPurple},
    mathescape=true,
	basicstyle=\scriptsize,
    keywords={input,output,begin,end,if,then,else,while,do,for,each,return}
	%tabsize=4,
	%breakatwhitespace=false,
	%breaklines=true,
	%captionpos=b,
	%keepspaces=true,
	%numbersep=5pt,
	%showspaces=false,
	%showstringspaces=false,
	%showtabs=false,
	%language=C
}

\usepackage{minted}


\usetheme[titleformat=smallcaps, numbering=fraction, background=light, progressbar=frametitle]{metropolis}

\title{Digital Technology}
\subtitle{Json CSV Pandas}

\author{Stefano Cereda\\
		stefano.cereda@polimi.it
	}
\date{21/04/2020}
\institute[PoliMi]{Politecnico Milano}
\logo{\includegraphics[width=15mm]{../logopolimi}}

\setbeamercovered{invisible}

\makeindex

\begin{document}
\begin{frame}
	\maketitle
\end{frame}

\begin{frame}{JSON}
    You have seen at lesson that JSON is built with name: value pairs.
    A JSON file can thus be described with a Python dictionary.

    The JSON module allows to convert a between python objects and their textual JSON representation.
\end{frame}

\begin{frame}[fragile]{JSON --- dumps}
    The \emph{dumps} function converts python variables to strings:
    \begin{minted}{Python}
import json
my_dict = {'firstName': 'John',
           'age': 25
          }
print(type(my_dict))

js_repr = json.dumps(my_dict)
print(js_repr)
print(type(js_repr))
    \end{minted}
\end{frame}

\begin{frame}[fragile]{JSON --- loads}
    The \emph{loads} function is used to convert from string to python objects:
    \begin{minted}{Python}
# continues from previous slide
my_obj = json.loads(js_repr)
print(my_obj)
print(type(my_obj))
    \end{minted}
\end{frame}


\begin{frame}[fragile]{JSON --- dump and load}
    Usually we want to read and write JSON \emph{files} instead of dealing with strings.

    The \emph{dump} and \emph{load} allow to work with files (no trailing \emph{s} as they do not work with
    \emph{s}trings.

    To use them, we first need to be able to open files in Python:
    \begin{minted}[fontsize=\tiny]{Python}
import json
my_dict = {'firstName': 'John', 'age': 25 }

# Open a file in write mode
with open('./my_file.json', 'w') as outfile:
    # outfile is a variable pointing to my_file.json, opened in write mode
    json.dump(my_dict, outfile)
# outfile is closed when we exit the with statement

# now we open the file in read mode and read its content:
with open('./my_file.json', 'r') as infile:
    loaded_obj = json.load(infile)

if loaded_obj != my_dict:
    print("JSON does not work!")
    \end{minted}
\end{frame}

\begin{frame}[fragile]{The open function}
    The \emph{open} function opens a file and return the corresponding file object.
    Its two most important parameters are the \emph{filename}, representing the path to the file to be opened, and the
    \emph{mode}: a string representing how we should open the file ('w' for write, 'r' for read).

    \url{https://docs.python.org/3/library/functions.html#open}
\end{frame}

\begin{frame}[fragile]{The with statement}
    When we finish processing the file we should close it:
    \begin{minted}{Python}
outfile = open('./my_file.json', 'w')
json.dump(my_dict, outfile)
outfile.close()
    \end{minted}
    However, doing this we must remember to close the file, even if the write instruction fails.

    The with statement is a much more convenient way to deal with file objects, as we have seen in previous examples.
\end{frame}



\begin{frame}[fragile]{Installing external modules with pip}
    We have seen how to import a module from the Python standard library:
    \begin{minted}{Python}
import math
from random import uniform
from json import dumps
    \end{minted}

    However, there exists many modules which are not included in the standard library (numpy, pandas, scipy,
    scikit-learn, etc\ldots)

    You can install, upgrade, and remove packages using a program called \emph{pip}:
    \begin{minted}{text}
pip install numpy
pip install pandas
pip uninstall numpy
pip freeze
    \end{minted}
\end{frame}

\begin{frame}[fragile]{Virtual environments}
    Different python applications will require different versions of the packages.

    The solution is to create a virtual environment, a directory that contains a
    Python installation plus a number of additional packages.

    Virtualenv creation:
    \begin{minted}{bash}
python -m venv my_venv
# equivalent: virtualenv my_venv
    \end{minted}

    And activation:
    \begin{minted}{text}
my_venv\Scripts\activate.bat  # Windows
source my_venv/bin/activate   # Linux and MacOS
    \end{minted}

    Using \emph{pip} with an active virtual environment will work on the virtual environment.
\end{frame}

\begin{frame}{PyCharm, pipenv and anaconda}
    In PyCharm every project has its own virtual environment, and you can easily manage its packages
    (File/Settings/Project/Project Interpreter).

    Other widely used programs for managing virtual environments are pipenv and anaconda.
\end{frame}

\end{document}
